\section{Conclusion}
\label{sec:conclusion}

This paper presented an OOB-guided cross-attentive RF-HSTU hybrid model for robust LoRa device authentication. The proposed architecture combines CNN-stem RF tokenization, RF patch sequence modeling, chirp-aware positional embeddings, and OOB-guided cross-attention with gated residual injection. The key idea is to use OOB spectral distortion as auxiliary hardware evidence while injecting it selectively into the main RF token sequence, rather than relying on naive feature concatenation.

On the public OSU LoRa dataset, under a strict source-domain validation cross-day protocol, the proposed model achieved 75.0$\pm$5.3\% file-level accuracy, compared with 54.2$\pm$14.2\% for the reproduced CNN-IQ baseline. The paired bootstrap confidence interval for the accuracy gain was entirely positive, supporting the stability of the improvement. Controlled fusion/chirp ablations further showed that cross-attentive OOB fusion is the primary mechanism behind the cross-day gain: removing chirp embedding from the cross-attentive model did not degrade accuracy, whereas adding chirp embedding to concat OOB fusion did not prevent collapse.

Extended deployment evaluations showed that the proposed hybrid model provides useful robustness under distance and indoor-location shifts, and the edge benchmark indicated that its inference overhead remains practical for gateway-side authentication. However, strict source-only cross-receiver transfer remained close to chance and direction-dependent. Receiver calibration mismatch therefore remains a major open challenge for deployable LoRa RFFI. A natural next step is receiver-adaptive calibration and multi-receiver training without target-receiver labels, together with open-set authentication and lightweight compression for edge gateways.
